\section{Maintenance-trigger errors and cost sensitivity}\label{sec:decision}
\subsection{Both errors, with engine-level uncertainty}
Table~\ref{tab:dec} and Figure~\ref{fig:dec}a,b give the heteroscedastic head at $L=20$. Under the sign-coherent positive drift on six-condition sub-datasets, premature triggers reach \JDriftPremRange\% of evaluated engines. Engine-level 95\% intervals span \JDriftPremCI\% over the four cells (half-width at most \JDriftPremHalfCI{} points), and the model-to-model range is \JDriftPremModelRange\%. No unrecognised need occurs: the rate is at most \JDriftUnrecMax\% in every one of the five mechanisms, eight backbone--sub-dataset pairs and three lead times. Premature rates under drift range \JDriftPremAllMech\% across the five mechanisms, so the choice of mechanism does not remove the error. The prematurely triggered engines have a mean unclipped true remaining life of \JDriftRULunc{} cycles (\JDriftRULcap{} clipped), far from needing maintenance. On single-condition sub-datasets drift triggers prematurely on \JDriftPremSingle\% (FD001, LSTM and Transformer) and \JDriftPremSingleThree\% (FD003) of engines.

\begin{table}[!t]
\caption{Terminal-instant decision stress test, heteroscedastic head, six-condition sub-datasets, $L=20$, one-sided 95\% lower bound. Rates are percentages of evaluated engines with engine-bootstrap 95\% intervals in brackets (2000 resamples, conditional on the five trained models); ``models'' is the range of the premature rate over the five models; ``cond.'' is the unrecognised rate among engines with true RUL $\le L$. Remaining life is the mean unclipped true RUL at the evaluated instant among prematurely triggered engines. All mechanisms and lead times: Supplement S4.}
\label{tab:dec}
\centering\footnotesize
\begin{tabular}{llcccccc}
\toprule
Cell & Condition & Premature (\%) & Models & Unrecognised (\%) & Cond.\ (\%) & Remaining life \\
\midrule
\tabin{tables/jrr_tab_decision.tex}\bottomrule
\end{tabular}
\end{table}

Random-sign bias errs on both sides, at lower rates than the positive drift. Premature triggers are \JBiasPremRange\% (all mechanisms \JBiasPremAllMech\%) and unrecognised needs \JBiasUnrecRange\% of evaluated engines, equivalent to \JBiasUnrecCondRange\% of the engines with true RUL $\le L$ (all mechanisms \JBiasUnrecAllMech\%). For the heteroscedastic head on the six-condition cells, clean inputs give \JCleanPremRange\% premature triggers and at most \JCleanUnrecMax\% unrecognised needs; noise at 1\% of range gives \JNoisePremRange\% premature triggers. The side of the error therefore depends on the sign pattern of the shift (Section~\ref{sec:coherent}): a decision audit has to count both sides, and a coverage figure does not say which side fails. The premature rates under drift vary with the lead time (\JDriftPremLten\% at $L=10$, \JDriftPremLthirty\% at $L=30$ over the four six-condition cells; Figure~\ref{fig:dec}c).

\begin{figure}[!t]
\centering
\includegraphics[width=\textwidth]{figures/fig3.pdf}
\caption{(a, b) Premature triggers and unrecognised needs of the heteroscedastic head at $L=20$ for each sub-dataset and condition; circles: LSTM, squares: Transformer; bars are engine-bootstrap 95\% intervals. (c) Lead-time dependence on the four six-condition cells; thin lines are individual cells, thick lines their mean. (d) Number of the 24 backbone--sub-dataset--lead-time cases per condition in which the ranking of the five mechanisms (solid) or only its top place (hatched) differs between cost ratios 5, 20 and 100.}
\label{fig:dec}
\end{figure}

\subsection{Mean or scale: an output swap at the decision point}\label{sec:swap}
A contracted scale raises the lower bound, so scale contraction alone cannot trigger earlier; an expanded scale can. Mean and scale of the heteroscedastic head were each taken from the clean or the perturbed condition of the same model. On six-condition sub-datasets the mean under positive drift falls by $\JMuShiftDrift$ cycles on average; changing only the mean raises premature triggers by \JSwapMuDrift{} points (lower interval limits at least \JSwapMuDriftLoCI{}), whereas changing the scale given the perturbed mean moves them by $\JSwapSigDrift$ points, at most \JSigExtMaxPt{} (\JSigExtMaxCell; interval \JSigExtMaxLo{} to \JSigExtMaxHi). On single-condition sub-datasets the scale contracts and offsets part of the mean effect in \JSwapSigDriftSingleN{} of four cases ($\JSwapSigDriftSingle$ points against a mean effect of \JSwapMuDriftSingle{}; largest \JSigExtMinPt{} for \JSigExtMinCell, interval \JSigExtMinLo{} to \JSigExtMinHi), which keeps premature rates low there. The scale-effect interval excludes zero in \JSigExtExclZero{} of \JSigExtCells{} cells, so the scale is not negligible. For random-sign bias the mean adds \JSwapMuBias{} points of premature triggers and \JSwapMuBiasUnrec{} points of unrecognised needs; the scale alone moves the unrecognised rate by $\JSwapSigBiasUnrec$ points. For the constant offsets of Section~\ref{sec:coherent} (Supplement S5), on six-condition data the mean effect is \JCPswapMuPrem{} points of premature triggers for $+5\%$ (scale effect \JCPswapSigPrem) and \JCNswapMuUnrec{} points of unrecognised needs for $-5\%$ (scale effect \JCNswapSigUnrec). The swaps identify the mean shift as the main driver, with scale changes of \JSigExtMinPt{} to \JSigExtMaxPt{} points at the extremes. This is an output-level decomposition, not a physical mechanism.

\subsection{Cost ratios and mechanism rankings}\label{sec:cost}
Across \JRevGroups{} backbone--sub-dataset--condition--lead-time groups, the cost ranking of the five mechanisms changes between cost ratios 5, 20 and 100 in \JRevTotal{} and the top-ranked mechanism changes in \JTopTotal. By condition (Figure~\ref{fig:dec}d), the full ranking changes in \JRevBias{} of \JNBias{} bias cases, \JRevNoise{} noise, \JRevClean{} clean and \JRevDrift{} drift cases; the top place changes in \JTopBias, \JTopNoise, \JTopClean{} and \JTopDrift{} cases respectively. Under drift the unrecognised rate is zero, so the cost is the same at every ratio and the ranking is the ranking of premature rates. Where the top place changes, the engine-bootstrap probability that the top mechanism at ratio 5 differs from that at ratio 100 has median \JTopBootMed{} (interquartile range \JTopBootIQR; at least 0.5 in \JTopBootGeFive of \JTopBootN{} cases). This describes how often resampling engines preserves the observed change of top place, conditional on the five models; it is not a significance test, and the probabilities for all 96 settings are in the supplementary CSV. The released per-engine bootstrap of the top mechanism against the other four excludes zero in only \NBootSig{} of 144 comparisons per backbone--sub-dataset. No mechanism is recommended on cost alone.

\subsection{Sensitivity checks}\label{sec:sens}
Perturbing the three operating-setting channels in addition to the sensors raised the premature rate under drift on FD004 from \NAOneFDfourSens{} (sensors only) to \NAOneFDfourJoint, while perturbing the settings alone gave \NAOneFDfourSet{}; on FD002 the two are indistinguishable (\NAOneFDtwoSens{} against \NAOneFDtwoJoint). The sensor-only premature rate is therefore the lower of the two in this FD004 check; this tested sensitivity does not make the sensor-only figures conservative in general. Engines with fewer than 30 cycles are excluded (\JExclTwo{} of 259 for FD002, \JExclFour{} of 248 for FD004); the effect of including them with padding was not assessed. Label-instant sensitivity is at most \JLabelHeadline{} points (Section~\ref{sec:framework}).
